\subsection{向量空间的对偶}

定理4. 对偶空间 \(\mathbf{E}^*\) （对应向量空间 \(\mathbf{E}\) ）的维数至少等于 \(\mathbf{E}\) 的维数。要使 \(\mathbf{E}^*\) 为有限维，必要且充分条件是 \(\mathbf{E}\) 为有限维，且此时 \(\dim \mathbf{E}^* = \dim \mathbf{E}\) .

如果 K 是 E 的标量域，则 E 同构于某个空间 \(\mathbf{K}_{s}^{\mathrm{(I)}}\) ，从而 \(E^{*}\) 同构于 \(K_{d}^{I}\) （§2，第6号，命题10）。由于 \(\mathbf{K}_{s}^{\mathrm{(I)}}\) 是 \(K_{d}^{I}\) 的一个子空间，故 \(\dim E = \text{Card}(I) \leqslant \dim E^{*}\) （第3号，命题4的推论4）；进一步，若 I 是有限的，则 \(\mathbf{K}_{d}^{I} = \mathbf{K}_{d}^{\mathrm{(I)}}\) （参见习题3(d)）。

推论。对于向量空间 \(\mathbf{E}\) ，关系 \(\mathbf{E} = \{0\}\) 和 \(\mathbf{E}^{*} = \{0\}\) 是等价的。

定理5. 给定向量空间（在同一域 K 上）的两个正合序列及线性映射

\[
0 \rightarrow \mathrm{E} ^ {\prime} \rightarrow \mathrm{E} \rightarrow \mathrm{E} ^ {\prime \prime} \rightarrow 0
\]

\[
0 \to \mathbf {F} ^ {\prime} \to \mathbf {F} \to \mathbf {F} ^ {\prime \prime} \to 0
\]

以及 K 上的两个向量空间 G、H，则相应的序列

\[
0 \rightarrow \operatorname{Hom} (\mathrm{E} ^ {\prime \prime}, \mathrm{G}) \rightarrow \operatorname{Hom} (\mathrm{E}, \mathrm{G}) \rightarrow \operatorname{Hom} (\mathrm{E} ^ {\prime}, \mathrm{G}) \rightarrow 0
\]

\[
0 \rightarrow \operatorname{Hom} (\mathrm{H}, \mathrm{F} ^ {\prime}) \rightarrow \operatorname{Hom} (\mathrm{H}, \mathrm{F}) \rightarrow \operatorname{Hom} (\mathrm{H}, \mathrm{F} ^ {\prime \prime}) \rightarrow 0
\]

是正合的且分裂的。

这由每个向量子空间都是直因子（第3号，命题4）这一事实以及§2第1号命题1和命题2得出。

推论。对于每个正合序列

\[
0 \longrightarrow \mathrm{E} ^ {\prime} \stackrel {u} {\longrightarrow} \mathrm{E} \stackrel {v} {\longrightarrow} \mathrm{E} ^ {\prime \prime} \longrightarrow 0
\]

（即同一域 K 上向量空间与线性映射的），则序列

\[
0 \longrightarrow \mathrm{E} ^ {\prime \prime *} \stackrel {t _ {v}} {\longrightarrow} \mathrm{E} ^ {*} \stackrel {t _ {u}} {\longrightarrow} \mathrm{E} ^ {\prime *} \longrightarrow 0
\]

是正合的且分裂的。

特别地，由此可知，对于 E 的每个向量子空间 M，典范同态 \(E^{*}/M^{\prime} \rightarrow M^{*}\) （其中 \(M^{\prime}\) 是 \(E^{*}\) 中与 M 正交的子空间（ \(\S 2\) ，第 4 号））是双射的。

定理 6. 对于域 K 上的每个向量空间 E，典范映射 \(c_{\mathbf{E}}: \mathbf{E} \to \mathbf{E}^{**}\) （§2，第 7 号）是单射的；要使它为双射，当且仅当 E 是有限维的。

第一个断言以及"若 E 是有限维的则 \(c_{E}\) 是双射"这一事实，都是§2第7号命题14的特殊情形。设 E 是无限维的，于是我们可以假定 \(\mathbf{E} = \mathbf{K}_{s}^{(L)}\) ，其中 L 是无限集，从而 \(E^{*} = K_{d}^{L}\) 。设 \((e_{\lambda})_{\lambda \in L}\) 是 E 的典范基，并令 \((e_{\lambda}^{*})_{\lambda \in L}\) 为 \(E^{*}\) 中相应的坐标线性型族（§2，第6号）； \(E^{*}\) 中由诸 \(e_{\lambda}^{*}\) 生成的向量子空间正是直和 \(F' = K_{d}^{(L)}\) ，而 L 为无限的假设蕴含 \(F' \neq E^{*}\) 。于是存在一个超平面 \(H'\) 于 \(E^{*}\) 中，包含 \(F'\) （第3号，命题8），且由于 \(E^{*}/H'\) 非零，其对偶亦非零（定理4的推论），该对偶等同于 \(H''\) ，即 \(H'\) 在 \(E^{**}\) 中的正交 （§2，第6号，命题9的推论）。但 \(H'' \cap c_{E}(E)\) 包含于按 \(c_{E}\) 取的像，即 \(F'\) 在 E 中的正交的像，而它按定义为 0；因此 \(c_{E}(E) = E^{**}\) 是不可能的。

通常将把 E 等同于 \(E^{**}\) 中 \(c_{E}\) 的像所成的子空间.

设 \(\mathbf{E}, \mathbf{F}\) 是域 \(\mathbf{K}\) 上的两个向量空间， \(u: \mathbf{E} \to \mathbf{F}\) 是一个线性映射。我们将定义典范同构：

\begin{enumerate}
\def\labelenumi{(\arabic{enumi})}
\item
  从 \(\operatorname{Im}(u) = u(\mathbf{E})\) 的对偶到 \(\operatorname{Im}(^t u) = {}^t u(\mathbf{F}^*)\) 上的典范同构。
\item
  从 \(\operatorname{Ker}(u) = \overline{u}^1(0)\) 的对偶到 \(\operatorname{Coker}(^t u) = \mathrm{E}^* / {}^t u(\mathrm{F}^*)\) 上的典范同构。
\item
  从 \(\operatorname{Coker}(u) = \mathbf{F} / u(\mathbf{E})\) 的对偶到 \(\operatorname{Ker}(^t u) = {}^t u^{-1}(0)\) 上的典范同构。
\end{enumerate}

我们写 \(\mathbf{I} = \operatorname{Im}(u),\mathbf{N} = \operatorname{Ker}(u),\mathbf{C} = \operatorname{Coker}(u)\) ；由正合序列

\[
0 \longrightarrow \mathrm{N} \longrightarrow \mathrm{E} \stackrel {p} {\longrightarrow} \mathrm{I} \longrightarrow 0, \quad 0 \longrightarrow \mathrm{I} \stackrel {j} {\longrightarrow} \mathrm{F} \longrightarrow \mathrm{C} \longrightarrow 0 \tag {14}
\]

我们通过转置（定理5的推论）推导出正合序列

\[
\begin{array}{l} 0 \longrightarrow \mathrm{I} ^ {*} \stackrel {{t _ {p}}} {{\longrightarrow}} \mathrm{E} ^ {*} \longrightarrow \mathrm{N} ^ {*} \longrightarrow 0, \\ 0 \longrightarrow \mathrm{C} ^ {*} \longrightarrow \mathrm{F} ^ {*} \stackrel {{t _ {j}}} {{\longrightarrow}} \mathrm{I} ^ {*} \longrightarrow 0. \end{array}\tag{15}
\]

此外，由于 \(u = j \circ p\) , \(^{t}u = ^{t}p \circ ^{t}j\) ；因此正合序列（15）定义了 \(C^*\) 到 \(\operatorname{Ker}(^{t}u)\) 上、 \(I^*\) 到 \(\operatorname{Im}(^{t}u)\) 上以及 \(N^*\) 到 \(\operatorname{Coker}(^{t}u)\) 上的典范同构，因为 \(^{t}p\) 是单射且 \(^{t}j\) 是满射。确切地说，设 \(y \in \operatorname{Im}(u)\) , \(z \in \operatorname{Ker}(u)\) , \(t \in \operatorname{Coker}(u)\) , \(y' \in \operatorname{Im}(^{t}u)\) , \(z' \in \operatorname{Coker}(^{t}u)\) , \(t' \in \operatorname{Ker}(^{t}u)\) ；当 \(y'\) , \(z'\) , \(t'\) 分别典范地等同于 \(\operatorname{Im}(u)\) , \(\operatorname{Ker}(u)\) 和 \(\operatorname{Coker}(u)\) 上的线性形式时，则

\[
\langle y, y ^ {\prime} \rangle = \langle x, y ^ {\prime} \rangle \quad \text { for   all } x \in \mathrm{E} \text { such   that } u (x) = y;\tag{16}
\]

\begin{enumerate}
\def\labelenumi{(\arabic{enumi})}
\setcounter{enumi}{16}
\item
  \(\langle z, z' \rangle = \langle z, x^* \rangle\) ，其中 \(x^* \in \mathbf{E}^*\) 取遍所有模 \(^t u(\mathbf{F})\) 的类等于 \(z'\) 的元素;
\item
  \(\langle t, t' \rangle = \langle s, t' \rangle\) ，其中 \(s \in \mathbf{F}\) 取遍所有模 \(u(\mathbf{E})\) 的类等于 \(t\) 的元素.
\end{enumerate}

特别地，我们从这些结果推导出：

命题10. 设E、F是同一域K上的两个向量空间，且 \(u: \mathbf{E} \to \mathbf{F}\) 是一个线性映射。

\begin{enumerate}
\def\labelenumi{(\roman{enumi})}
\item
  对于 \(u\) 为单射（相应地，满射），其必要且充分条件是 \(^{t}u\) 是满射（相应地，单射）。
\item
  \(\operatorname{rg}(u) \leqslant \operatorname{rg}(^t u)\) ，并且关系 \(\operatorname{rg}(u) = \operatorname{rg}(^t u)\) 在 \(\operatorname{rg}(u)\) 有限时成立。
\end{enumerate}

第二个断言由上述内容及定理4得出。

定理7. 设E是域K上的向量空间，F是E的子空间，F'是F在E*中的正交补。

\begin{enumerate}
\def\labelenumi{(\roman{enumi})}
\item
  \(\dim \mathbf{F}' \geqslant \operatorname{codim}_{\mathbb{E}} \mathbf{F}\) ；对于 \(\dim \mathbf{F}'\) 为有限，必要且充分条件是 \(\operatorname{codim}_{\mathbb{E}} \mathbf{F}\) 有限，且此时 \(\dim \mathbf{F}' = \operatorname{codim}_{\mathbb{E}} \mathbf{F}\) .
\item
  \(\mathbf{F}'\) 在 \(\mathbf{E}\) 中的正交补等于 \(\mathbf{F}\) .
\item
  每个有限维子空间 \(\mathbf{G}'\) （ \(\mathbf{E}^*\) 的）都是 \(\mathbf{E}\) 的某个子空间的正交补，该子空间必然等于 \(\mathbf{G}'\) 在 \(\mathbf{E}\) 中的正交补，且具有有限余维数。
\item
  我们知道 \(\mathbf{F}'\) 同构于对偶 \((\mathbf{E} / \mathbf{F})^*\) （§ 2，第 6 号，命题 9 的推论），因此该断言由定理 4 得出，因为 \(\dim (\mathbf{E} / \mathbf{F}) = \operatorname{codim}_{\mathbf{E}}\mathbf{F}\) 根据定义。
\item
  设 \(F_{1}\) 是 \(F'\) 在 E 中的正交补；显然 \(F \subset F_{1}\) ，且 \(F_{1}'\) （ \(F_{1}\) 的正交补）等于 \(F'\) （§2，第4号）；典范线性映射 \((\mathrm{E}/\mathrm{F}_{1})^{*} \to (\mathrm{E}/\mathrm{F})^{*}\) （即 \(E/F \to E/F_{1}\) 的转置）因此是双射的（§2，第6号，命题10的推论）；于是由命题10可知，典范映射 \(E/F \to E/F_{1}\) 是双射的，这蕴含 \(F_{1} = F\) .
\item
  设 \(\mathbf{G}'\) 是 \(\mathbf{E}^*\) 的一个维数为有限数 \(p\) 的子空间，并设 \(\mathbf{F}\) 为它在 \(\mathbf{E}\) 中的正交补；则 \(\operatorname{codim}_{\mathbf{E}} \mathbf{F} \leqslant \dim \mathbf{G}'\) 。事实上，若 \((a_{i}^{*})_{1 \leqslant i \leqslant p}\) 是 \(\mathbf{G}'\) 的一组基，则 \(\mathbf{F}\) 是线性映射 \(x \mapsto (\langle x, a_{i}^{*} \rangle)\) （从 \(\mathbf{E}\) 到 \(\mathbf{K}_{s}^{p}\) ）的核，其秩至多为 \(p\) （第4号，命题9），由此得出结论（第4号）。然后设 \(\mathbf{F}'\) 是 \(\mathbf{F}\) 在 \(\mathbf{E}^*\) 中的正交补；由（i）可知 \(\dim \mathbf{F}' \leqslant \dim \mathbf{G}'\) ；但另一方面显然有 \(\mathbf{G}' \subset \mathbf{F}'\) ，因此 \(\mathbf{F}' = \mathbf{G}'\) ( \(\S 2\) ，第3号，命题4的推论4）。
\end{enumerate}

注记。一个无限维子空间 \(G'\) （含于 \(E^{*}\) ）不一定是 E 的某个子空间的正交补，换句话说，若 F 是 \(G'\) 在 E 中的正交补，则 F 的正交补 \(F'\) （在 \(E^{*}\) 中）可能不同于 \(G'\) （练习20(b)）。†

推论1. 设 \((x_{i}^{*})_{1\leqslant i\leqslant p}\) 是 E 上的线性型的一个有限序列，并设 F 为 E 中由使得下式成立的 \(x\) 所组成的子空间，

\[
\langle x, x _ {i} ^ {*} \rangle = 0 \quad for 1 \leqslant i \leqslant p.
\]

那么 \(\operatorname{codim}_{E} F\) 等于诸 \(x_{i}^{*}\) 之集的秩，并且 E 上每个在 F 上为零的线性型都是诸 \(x_{i}^{*}\) 的线性组合。进而 \(\operatorname{codim}_{E} F \leqslant p\) ，并且为使 \(\operatorname{codim}_{E} F = p\) ，必要且充分条件是诸 \(x_{i}^{*}\) 线性无关。

集合 \(G'\) （由诸 \(x_{i}^{*}\) 的线性组合组成）是 \(E^{*}\) 的子空间，且 F 是 \(G'\) 在 E 中的正交补，因此根据定理7， \(\operatorname{codim}_{E} F = \dim G'\) ；进一步 \(\dim G' \leqslant p\) ，而关系 \(\dim G' = p\) 意味着 \((x_{i}^{*})\) 是一个自由系（第2号，命题1）；由此得到该推论。

推论2. (i) 设 \((x_{i}^{*})_{1\leqslant i\leqslant p}\) 是 E 上的线性型的一个有限序列。要使 \((x_{i}^{*})\) 成为一个自由系，必要且充分条件是存在 E 的元素的一个序列 \((x_{i})_{1\leqslant i\leqslant p}\) 使得 \(\langle x_i,x_j^*\rangle = \delta_{ij}\) （克罗内克指标）。

\begin{enumerate}
\def\labelenumi{(\roman{enumi})}
\setcounter{enumi}{1}
\tightlist
\item
  设 \((x_{i})_{1\leqslant i\leqslant p}\) 是 E 的元素的一个有限序列。要使 \((x_{i})\) 成为一个自由系，必要且充分条件是存在 E 上的线性型的一个序列 \((x_{i}^{*})_{1\leqslant i\leqslant p}\) ，使得 \(\langle x_i,x_j^*\rangle = \delta_{ij}\) .
\end{enumerate}

显然，(ii) 由 (i) 得出，只需将 E 视为 \(\mathbf{E}^{**}\) 的一个子空间（借助 \(c_{\mathbf{E}}\) ，定理 6）。设 \(\mathbf{G}'\) 为 \(\mathbf{E}^*\) 中由诸 \(x_{i}^{*}\) 生成的子空间，而 \(\mathbf{F}\) 为它在 \(\mathbf{E}\) 中的正交补； \(\mathbf{E} / \mathbf{F}\) 和 \(\mathbf{G}'\) 中的每一个都可在典范意义下等同于另一个的对偶；若族 \((x_{i}^{*})\) 是自由的，则在 \(\mathbf{E} / \mathbf{F}\) 中存在一组基 \((\dot{x}_i)\) ，它是 \((x_{i}^{*})\) 的对偶基，而 \((x_{i})\) 作为诸类 \(\dot{x}_{i}\) 的任意代表系，都具有所需的性质。反之，系统 \((x_{i})\) （满足 \(\langle x_{i}, x_{j}^{*} \rangle = \delta_{ij}\) ）的存在意味着：对所有 \(i\) ， \(\mathbf{E}^*\) 中与 \(\mathbf{K}\) . \(x_{i}\) 正交的子空间包含诸 \(x_{j}^{*}\) （其指标 \(j \neq i\) ） ，但不包含 \(x_{i}^{*}\) ，因此系统 \((x_{i}^{*})_{1 \leqslant i \leqslant p}\) 是自由的。

推论3. 设 S 是一个集合，V 是由 S 到 K 的映射组成的右向量 K-空间 \(K_{d}^{S}\) 的一个向量子空间。要使 \(\dim V \geqslant p\) （其中 p 是整数），必要且充分条件是存在 S 的 p 个元素 \(s_{i}\) 以及 V 的 p 个元素 \(f_{i}\) （ \(1 \leqslant i \leqslant p\) ），使得 \(f_{i}(s_{j}) = \delta_{ij}\) .

该空间 \(K_{d}^{S}\) 典范地等同于 \(E = K_{d}^{(S)}\) 的对偶空间，且 \(f(s) = \langle e_{s}, f \rangle\) 对 \(s \in S\) 和 \(f \in K_{s}^{S}\) 成立 ，其中 \((e_{s})_{s \in S}\) 是 E 的典范基。推论2随后表明该条件是充分的。反之，假设 \(\dim V \geqslant p\) ，因此存在 V 的一个维数为 p 的子空间 \(G'\) ；设 F 为 \(G'\) 在 E 中的正交补，从而 \(\dim(E/F) = p\) 。由第 1 号定理 2 可知，存在 p 个元素 \(s_{i} \in S\) 使得 \(e_{s_{i}} (1 \leqslant i \leqslant p)\) 构成 F 在 E 中的一个补空间的一组基（将定理 2（第 1 号）应用于一个由 E 的一组基组成的自由子集，以及由此自由子集与 E 的典范基的并集组成的生成系）；然后我们取 \(f_{i}\) 为 G' 的一组基的元素，该基对偶于由诸 \(e_{s_{i}} \mod F\) 的类组成的 E/F 的基.

推论4. 设E是一个向量空间，M、N是E的两个有限余维子空间；若M'、N'分别是M和N在E*中的正交补，则M ∩ N在E*中的正交补为M' + N'。

由于 M（相应地，N）是 \(M'\) （相应地 \(N'\) ）在 E 中的正交补（定理 7）， \(M \cap N\) 是 \(M' + N'\) 在 E 中的正交补，因此 \(M' + N'\) 是 \(M \cap N\) 在 \(E^*\) 中的正交补（定理 7 (iii)）。

推论 5. 设 E 是维数为 n 的有限维向量空间。～对于 E 的每个维数为 p 的子空间 F，F 在 E* 中的正交补 F' 的维数为 n - p。～对于 E* 的每个维数为 q 的子空间 G'，G' 在 E 中的正交补 G 的维数为 n - q，且 G' 是 G 在 E* 中的正交补。

定理 7 给出了 E 中超平面的另一个刻画：

命题 11. 对于向量空间 E 中的每个超平面 H，存在 E 上的一个线性形式 \(x_0^*\) ，使得 \(\mathrm{H} = x_0^{*-1}(0)\) 。给定这样的一个形式 \(x_0^*\) ，对于 E 上的线性形式 \(x^*\) ，要满足 \(\mathrm{H} = x^{*-1}(0)\) ，必要且充分的条件是 \(x^* = x_0^*\alpha\) ，其中 \(\alpha\) 是一个标量 \(\neq 0\) 。反之，对于 E 上的每一个线性形式 \(x^* \neq 0\) ，子空间 \(x^{*-1}(0)\) 是 E 的一个超平面。

这个陈述仅仅表达了定理7对于E的余维数为1的子空间以及E*的维数为1的子空间的情形。

如果 H 是一个超平面，且 \(x_{0}^{*}\) 是满足 \(\mathrm{H}=x^{*-1}(0)\) 的一个线性型，则关系

\[
\langle x, x _ {0} ^ {*} \rangle = 0
\]

它刻画了元素 \(x \in \mathbf{H}\) ，称为 \(\mathbf{H}\) 的一个方程.

更一般地，如果 \((x_{\iota}^{*})\) 是 E 上的一族线性型，且 F 表示诸超平面 \(x_{\iota}^{*-1}(0)\) 的交所构成的向量子空间，则关系式“对所有 \(\iota\) , \(\langle x, x_{\iota}^{*}\rangle = 0\) 成立”刻画了 F 的元素；这些关系式

\[
\langle x, x _ {\iota} ^ {*} \rangle = 0 \quad \text { for   all } \iota
\]

构成子空间 F 的一个方程组。定理 7 (ii) 表达了 E 的每个向量子空间都可以由一个方程组来定义这一事实。

定理 7 (i) 和 (ii) 还证明了，余维数为有限数 p 的子空间 F 可以由 p 个方程组成的方程组

\[
\langle x, x _ {i} ^ {*} \rangle = 0, \quad 1 \leqslant i \leqslant p,\tag{19}
\]

来定义，其中这些型 \(x_{i}^{*}\) 是线性无关的。反之，定理 7 的推论 1 表明，子空间 \(F\) 若由 \(p\) 个方程 (19) 组成的方程组所定义，则其余维数 \(\leqslant p\) ，且其余维数为 \(p\) 当且仅当诸 \(x_{i}^{*}\) 线性无关；这等价于说 F 不能由至多包含 (19) 中 p - 1 个方程的方程组来定义。

